\section{Introduction and current theorem structure}\label{sec:current-overview}\label{sec:introduction}
A finite physical system, a finite classical model of its observations, and a finite observer implementing an optimal decision are different mathematical objects. The central question of General Theta Foundations I is how a prepared causal experiment turns physically acquired information into attainable prediction under explicit resources. The actual-history transfer theorem compares task risk with acquired geometry under its paid converse and common-construction certificates. The performed-law theory establishes unrestricted attainment, and the boundary results determine when exact observer memory follows. We now supply a constructive reduction of a broad class of controlled physical instruments without conflating any of these implications.

Theorem~\ref{thm:r68-instrument} replaces a supplied invariant-cone certificate by a check on each actual command/report branch. If all nonzero normalized outputs have a common positive spectral floor, one pure-state net works for every branch, even when commands do not commute or share an invariant state. The resulting hidden classical instrument preserves the complete observations of every adaptive observer and every terminal measurement. Its expected-cost identity is proved at each actual begun history and summed monotonically; it is not inferred from a bounded-loss discrepancy. A block extension counts all intermediate reports and its clock rather than making block calls free.

Theorem~\ref{thm:r68-noise-exponent} gives matching powers of the mixture margin for uniform state-positive realization at each fixed physical dimension. The exponent comes from the geometry of pure-state projective caps. Mixed hidden density vertices cannot evade that bound. Finite physical command families already force the same exponent, although their command count grows. This is deliberately not a universal lower bound for all hidden response representations: a classical generator need not correspond to a quantum density matrix when only one task is prescribed. The separate fixed-angle spectral bound, Corollary~\ref{cor:r68-qubit-lower}, applies to all stochastic scalar-response realizations and is explicitly non-sharp.

Positive realization, invariant polyhedral cones and finite-rank obstructions are established ideas \cite{Heller65,CarlylePaz71,MonrasWinter14}. Fanizza, Lumbreras and Winter also obtain noise-dependent classical-memory obstructions for quantum processes \cite{Fanizza24}. The contribution asserted here is the explicit branchwise construction, its dimension-dependent uniform count and sparse rows, the exact adaptive stopped-resource identity, and the distinctions between its resource classes. We do not claim to originate the cone mechanism or the classical--quantum separation.

The previously proved single-command boundary criterion and squared-audit hierarchy remain in the main article below, along with the complete actual-history comparison proof. The robust optimal face includes global modelwise inequalities even when every Bayes action is tied. Neither a finite physical hidden model nor a finite linear Hankel rank alone settles the exact implementation of that face. The full performed-law, primal--dual, timing, finite-width, reset, spectral, geometry and predictive-quotient developments remain current in the paired supplement; all predecessor sources remain immutable ancestry.
